\documentclass[../main.tex]{subfiles}

\begin{document}
	
	\chapter{数学视角下的演化、结构与稳定性}
	
	\begin{motivation}
		你可能会问：为什么需要数学附录？正文不是已经说清楚了吗？
		
		\textbf{核心动机}：数学不是炫耀，而是让那些"大概如此"的洞察变得可检验、可复现、可争议。当你看到一个关于演化稳定性的论断时，附录会告诉你：这个结论依赖哪些假设？如果把这些假设放松，结论会如何变化？这就是形式化的力量。
	\end{motivation}
	
\begin{logicmap}
	正文直觉与洞察
	$\rightarrow$ 需要形式化以检验
	$\rightarrow$ 数学附录模块
	$\rightarrow$ 对应各章核心论断
	$\rightarrow$ 可复现、可争议、可验证
\end{logicmap}

	\section{V.0 使用说明与工程声明}
	
	\subsection{你有没有想过，为什么你需要阅读这个数学附录？}
	你可能正在想：正文已经用文字说清楚了，为什么还要看公式？这个附录不是又要讲一遍新东西吧？实际上，数学不是炫耀，而是让那些"大概如此"的洞察变得可检验、可复现、可争议。当你看到一个关于演化稳定性的论断时，附录会告诉你：这个结论依赖哪些假设？如果把这些假设放松，结论会如何变化？这就是形式化的力量。正确使用本附录的方法是按需跳转：不必通读，只在你想要验证某个直觉、或者发现矛盾时，打开对应的数学模块，检查它的公理边界。
	
	% 动机说明
	% 本节目的：防止读者误将数学附录视为新的叙事主线
	\begin{mathbox}{附录定位与阅读方式}
		\textbf{定位}：本附录不是新的论述层，而是对正文关键论断的\textbf{形式化、可核验重写}。
		
		\textbf{阅读建议}：
		\begin{itemize}
			\item 可在阅读正文后，按需跳转至对应数学模块；
			\item 每一节均通过超链接与正文一一对应；
			\item 附录中不引入新的经验性主张，仅展开假设、模型与边界。
		\end{itemize}
	\end{mathbox}
	
	\subsection{逻辑衔接}
	明确了附录的定位与使用方法后，我们不再引入新的叙事，而是直接进入正文关键论断的形式化展开。下一模块（chapter1_existence_dynamics）将探讨存在动力学中的演化稳定性条件，用博弈论的形式语言揭示动机系统的数学基础。
	
	% 核心数学模块
	\subfile{appendix/chapter1_existence_dynamics}
	\subfile{appendix/chapter2_structure_info}
	\subfile{appendix/chapter3_history_probability}
	\subfile{appendix/chapter4_selforg_dynamics}
	\subfile{appendix/chapter5_checklist}
	
	% 具体数学内容（原始模块）
	
\end{document}
